% Chapter 10 of Optimizers in the Wild: gradient clipping, generalized smoothness and loss spikes.
\newcommand{\spikescliff}{\Loss_{\mathrm{cliff}}}
\newcommand{\spikesplat}{\Loss_{\mathrm{plat}}}
\newcommand{\spikesclip}{\operatorname{clip}}

\casechapter{Clipping and Spikes}{Generalized smoothness and the anatomy of an instability}{Gradient clipping, z-loss, QK-norm}{ch:spikes}
\begin{refsection}[chapters/10-spikes.bib]

\begin{chaptersummary}
Large language models are trained with the norm of the gradient clipped at $1$, a penalty on the normalizer of the output softmax and normalized queries and keys, and their losses still spike. The idealization is $(L_0,L_1)$-smoothness, a bound on the curvature by $L_0$ plus $L_1$ times the norm of the gradient. We prove a sharp local smoothness lemma, show that clipped steepest descent reaches a gradient of size $\eps$ within $\max\{4L_0\Delta_0/\eps^2,6L_1\Delta_0/\eps\}$ steps while every fixed step size pays almost linearly for the largest gradient met on the way, and show that normalized updates are clipped by construction. We then take three instabilities apart exactly: a block of weights that wakes up after a quiet period receives Adam steps of up to $0.21/\sqrt{1-\beta_2}$ times the learning rate; the cross-entropy has a flat direction along which adaptive updates drift and which the z-loss removes; and attention logits grow with the square of the weight scale unless queries and keys are normalized. The theory explains why each defence helps; it does not predict when a spike will come, and real spikes have causes it leaves out.
\end{chaptersummary}

\section{In the wild: clipped gradients and spiking losses}\label{sec:spikes:wild}

Almost every large training run rescales its gradient before using it. The Euclidean norm of the gradient is computed over all parameters at once; if it exceeds a threshold, the whole gradient is multiplied by the threshold divided by the norm, and otherwise it is left alone. The rule comes from recurrent networks, whose losses have \term{cliffs}\index{cliff}: thin regions in which the gradient and the curvature are both enormous, so that a single gradient step throws the parameters far away.\sidenote{Clipping entered deep learning with recurrent language models \parencite{Mikolov2012}. \Textcite{PascanuMikolovBengio2013} proposed clipping the norm of the gradient and explained the need with the geometry of cliffs, building on the analysis of exploding and vanishing gradients by \textcite{BengioSimardFrasconi1994}.} Transformers inherited the rule, and with it a growing list of defences against the instabilities that clipping does not prevent.

\begin{practice}[GPT-3, \textcite{BrownEtAl2020}; LLaMA, \textcite{TouvronEtAl2023}; PaLM, \textcite{ChowdheryEtAl2023}; ViT-22B, \textcite{DehghaniEtAl2023}; Kimi K2, \textcite{KimiTeam2025}]
Clip the global norm of the gradient at $1.0$ before the optimizer step (GPT-3, LLaMA, PaLM). Use Adam with $\beta_2=0.95$ rather than the default $0.999$ (GPT-3, LLaMA). Add $10^{-4}(\log Z)^2$ to the cross-entropy, where $Z$ is the normalizer of the softmax over the vocabulary (the \term{z-loss}\index{z-loss} of PaLM). Apply a LayerNorm to the queries and the keys before their dot product (\term{QK-norm}\index{QK-norm}, ViT-22B). With Muon, after each step rescale the query and key weights of every attention head whose largest logit on the batch exceeds a threshold, $100$ for Kimi K2 (\term{MuonClip}\index{MuonClip}). When the loss spikes anyway, restart from a checkpoint shortly before the spike and skip the batches around it (PaLM).
\end{practice}

The values travel from recipe to recipe; we know of no derivation of the threshold $1.0$ or of the coefficient $10^{-4}$, and in that sense both are folklore. What has been observed is the following.

\begin{observation}[Loss spikes; \textcite{ChowdheryEtAl2023}]\label[observation]{obs:spikes:palm}
While the $540$-billion-parameter PaLM model was trained, with the gradient norm clipped at $1.0$, the loss spiked roughly $20$ times, at irregular intervals and sometimes late in training. Restarting from a checkpoint about $100$ steps before a spike and skipping about $200$ to $500$ batches avoided it. Replaying the batches around a spike from a different, earlier checkpoint did not produce a spike.
\end{observation}

The last sentence carries the weight: the spike was in neither the data nor the parameters alone, but in their combination. The report is qualitative and concerns one run. The OPT team describes the same experience \parencite{ZhangEtAl2022OPT}: their divergences were cured by lowering the learning rate and restarting from an earlier checkpoint, and they coincided with a collapse of the dynamic loss scale used for mixed precision and with spikes in the norm of the activations of the last layer.

\begin{observation}[Curvature grows with the gradient; \textcite{ZhangHeSraJadbabaie2020}]\label[observation]{obs:spikes:L0L1}
Along the training trajectory of an LSTM language model, a local smoothness estimated from the change of the gradient between nearby iterates was strongly correlated with the norm of the gradient: where the gradient was large, so was the curvature.
\end{observation}

If the observation holds, a global bound on the curvature is set by the steepest part of the trajectory and is far larger than the curvature met during most of the run, so the step size $1/L$ that the classical theory prescribes is wasted almost everywhere. The evidence is a correlation along trajectories of a few networks, not a law.

\begin{observation}[Two instabilities in small models; \textcite{WortsmanEtAl2024}]\label[observation]{obs:spikes:proxies}
Small transformers trained at high learning rates reproduce two instabilities first reported for large runs: the growth of the attention logits, which makes the attention weights nearly one-hot \parencite{DehghaniEtAl2023}, and the divergence of the output logits from log-probabilities \parencite{ChowdheryEtAl2023}. Normalizing queries and keys removes the first, and the z-loss removes the second.
\end{observation}

\begin{observation}[Spikes follow a stale second moment; \textcite{WortsmanEtAl2023}]\label[observation]{obs:spikes:stale}
In training large CLIP models with AdamW, loss spikes followed, within a few iterations, steps at which the squared gradients of some layers were much larger than AdamW's running estimate of them: the second-moment estimate was out of date.
\end{observation}

From runs of language models with $7$ to $546$ billion parameters, \textcite{MolybogEtAl2023} (unrefereed) argue for a related account: Adam can enter a state in which its update has a large norm and is nearly uncorrelated with the direction of descent, and this is more likely in deep models trained with large batches.

The reports share a pattern. Each instability is a step too long for the curvature where it lands, and in each the step became too long because some quantity grew that the loss did not register in time: the gradient near a cliff, the gap between Adam's estimate of the gradient's scale and the scale itself, a common shift of the output logits, the size of queries and keys. Why can a fixed step size not cope with curvature that grows with the gradient, and what exactly does clipping buy? Through which mechanisms does a well-tuned run produce a step that is far too long? The idealization of this chapter answers the first two questions with theorems and the third with three exact mechanisms. It does not predict when a spike happens.

\section{The idealization: curvature that grows with the gradient}\label{sec:spikes:ideal}

Classical smoothness bounds the curvature by one number everywhere. On a cliff that number is huge, and the step size it allows is wasted on the rest of the landscape. \Textcite{ZhangHeSraJadbabaie2020} let the bound grow with the gradient.\sidenote{They wrote the condition in the Euclidean norm. The general norm costs nothing in the proofs and is what \cref{prop:spikes:normalized} needs: sign descent and Muon are steepest descent in the $\ell_\infty$ and spectral norms (\cref{ch:sign,ch:spectral}).}

\begin{assumption}[$(L_0,L_1)$-smoothness]\label[assumption]{ass:spikes:L0L1}
Let $\norm{\cdot}$ be a norm on $\R^d$ with dual norm $\norm{\cdot}_*$, and for a symmetric matrix $H$ let $\norm{H}=\max_{\norm u\le1}\norm{Hu}_*$. The objective $\Loss\colon\R^d\to\R$ is twice continuously differentiable and bounded below, $\Loss^\star=\inf\Loss>-\infty$, and there are constants $L_0\ge0$ and $L_1>0$ such that
\begin{equation}\label{eq:spikes:L0L1}
  \norm{\nabla^2\Loss(\theta)}\le L_0+L_1\norm{\nabla\Loss(\theta)}_*\qquad\text{for every }\theta\in\R^d .
\end{equation}
\end{assumption}

\term{$(L_0,L_1)$-smoothness}\index{generalized smoothness} with $L_1=0$ would be ordinary $L_0$-smoothness; the point of $L_1>0$ is to allow curvature that grows without bound, provided the gradient grows with it. The exponential $e^x$ is $(0,1)$-smooth and $\cosh x-1$ is $(1,1)$-smooth; the logistic loss $\log(1+e^{-x})$ is both $(\frac14,0)$- and $(0,1)$-smooth; $x^4$ is $(L_0,L_1)$-smooth exactly when $L_0L_1^2\ge16$; and $\sum_ia_i(\cosh\theta_i-1)$ is $(\max_ia_i,1)$-smooth in the Euclidean norm and $(\sum_ia_i,1)$-smooth in $\ell_\infty$ (\cref{ex:spikes:constants}). A quadratic with Hessian $H$ is $(\norm H,0)$-smooth and no better, because at its minimizer the gradient vanishes. The function that saturates \eqref{eq:spikes:L0L1} everywhere is the \term{exponential cliff}\index{cliff!exponential}
\begin{equation}\label{eq:spikes:cliff}
  \spikescliff(x)=\frac{L_0}{L_1^2}\bigl(e^{L_1\abs x}-1-L_1\abs x\bigr),\qquad
  \abs{\spikescliff'(x)}=\frac{L_0}{L_1}\bigl(e^{L_1\abs x}-1\bigr),\qquad
  \spikescliff''(x)=L_0e^{L_1\abs x},
\end{equation}
a valley of curvature $L_0$ whose walls steepen exponentially; it is twice continuously differentiable, and $\spikescliff''=L_0+L_1\abs{\spikescliff'}$. Exploding gradients have this shape: in a recurrent network the loss depends on a recurrent weight through high powers of it \parencite{BengioSimardFrasconi1994,PascanuMikolovBengio2013}, and a polynomial of high degree has derivatives that grow together.

\paragraph{The method.} Write $g^\sharp=-\norm g_*\lmo(g)$ for the sharp operator of the norm (\cref{ch:steepest}; \cref{lem:toolkit:conjugate}), so that $\ip{g,g^\sharp}=\norm g_*^2$ and $\norm{g^\sharp}=\norm g_*$. \term{Clipped steepest descent}\index{gradient clipping} with step size $\eta$ and largest step length $\gamma$ is
\begin{equation}\label{eq:spikes:clipped}
  \theta_{t+1}=\theta_t-\eta_t\,g_t^\sharp,\qquad g_t=\nabla\Loss(\theta_t),\qquad \eta_t=\min\Bigl\{\eta,\frac{\gamma}{\norm{g_t}_*}\Bigr\},
\end{equation}
with $\eta_t=\eta$ when $g_t=0$. Its step has length $\min\{\eta\norm{g_t}_*,\gamma\}$. Below the threshold $\norm{g_t}_*=\gamma/\eta$ it is unnormalized steepest descent with step size $\eta$; above it, it is normalized steepest descent $\theta_{t+1}=\theta_t+\gamma\lmo(g_t)$ with step length $\gamma$. In the Euclidean norm $g^\sharp=g$, and \eqref{eq:spikes:clipped} is the rule of the practice box: $\theta_{t+1}=\theta_t-\eta\spikesclip_\tau(g_t)$ with $\spikesclip_\tau(g)=\min\{1,\tau/\norm g_2\}g$ and threshold $\tau=\gamma/\eta$. The norm is taken over all parameters of the network at once, as in PyTorch's \texttt{clip\_grad\_norm\_}; clipping each layer separately is clipped steepest descent in a product norm (\cref{ch:modular}).\sidenote{Clipping \emph{by value} instead clips each coordinate at the threshold; it is the trust-region step for a box rather than a ball (\cref{ex:spikes:trust}), and it changes the direction of the step.} Clipping is a trust region\index{trust region}: the clipped step minimizes $\ip{g,u}+\norm u_2^2/(2\eta)$ over the ball $\norm u_2\le\gamma$ (\cref{ex:spikes:trust}).

\begin{fieldguide}{Clipped steepest descent}{\textcite{PascanuMikolovBengio2013}, Euclidean norm}
\fgrow{update}{$\theta_{t+1}=\theta_t-\min\{\eta,\gamma/\norm{g_t}_*\}\,g_t^\sharp$; in $\ell_2$, $\theta_{t+1}=\theta_t-\eta\min\{1,\tau/\norm{g_t}_2\}\,g_t$, threshold $\tau=\gamma/\eta$}
\fgrow{norm}{every norm; the step is unnormalized below the threshold and normalized, of length $\gamma$, above it}
\fgrow{preconditioner}{none; in practice the clipped gradient is passed to Adam, whose moments see the clipped values}
\fgrow{averages}{none}
\fgrow{state}{none; one norm per step, over all parameters}
\fgrow{in the wild}{$\tau=1.0$ on the global $\ell_2$ norm, in front of AdamW (GPT-3, LLaMA, PaLM)}
\fgrow{variants}{clipping by value, an $\ell_\infty$ trust region (\cref{ex:spikes:trust}); clipping the update rather than the gradient \parencite{ShazeerStern2018}}
\end{fieldguide}

\begin{dictionary}
$(L_0,L_1)$-smoothness \eqref{eq:spikes:L0L1} & curvature that is largest where the gradient is largest: cliffs, the first steps of training, a spike\\
$L_0$ & the curvature of the flat part of the landscape, which sets the useful learning rate\\
$L_1\norm{\nabla\Loss}_*$ & the extra curvature near a cliff or during a spike\\
$M$, the largest gradient on the sublevel set & the largest gradient norm the run will meet, including its spikes\\
step size $\eta$; step length $\gamma$; $\gamma/\eta$ & learning rate; learning rate times threshold; the threshold \texttt{max\_norm}\\
a coordinate that is silent, then loud & a token's embedding, a head or an expert whose gradient was near zero and suddenly is not\\
$a_k$ of \eqref{eq:spikes:burst} & Adam's update, per coordinate and in units of the learning rate, $k$ steps after waking up\\
the flat direction $\mathbf1$ of the logits & a common shift of all output logits, which moves $\log Z$\\
$\alpha(\log Z)^2$ & the z-loss, $\alpha=10^{-4}$\\
$s=\ip{W_Qx,W_Kx'}/\sqrt{d_{\mathrm h}}$ & an attention logit; MuonClip watches its maximum per head\\
the gains $\omega_q,\omega_k$ & the scale parameters of the normalizations of queries and keys
\end{dictionary}

\paragraph{What was thrown away.} The gradient is exact in the main theorems; noise returns only in \cref{thm:spikes:heavy} and in the anatomy of \cref{sec:spikes:anatomy}. The network is reduced to one relation between its curvature and its gradient, with constants that do not change during training; the constants of a real run do change, and we cannot compute them. Finite precision is absent. These are the right first simplifications because the failure of fixed step sizes is already present without noise or architecture: a step size that is safe where the gradient is largest is too small everywhere else, and \cref{thm:spikes:lower} shows that this costs a factor that grows with the largest gradient.

\paragraph{What one had to notice.} Classical analyses take a global smoothness constant as given and derive a step size from it. \Textcite{ZhangHeSraJadbabaie2020} measured the curvature along trajectories and saw that it moved with the gradient norm. Once curvature is tied to the gradient, a step whose \emph{length} is bounded is safe everywhere, and bounding the length of the step is exactly what clipping does. For spikes, the thing to notice is that each defence in the practice box targets a quantity whose growth the loss does not see until it is too late: the optimizer's state, a flat direction of the loss, a scale.

\section{The theory: clipping under generalized smoothness}\label{sec:spikes:clipping}

Everything rests on one lemma: under \eqref{eq:spikes:L0L1}, the gradient cannot change faster than exponentially along a segment, at a rate set by $L_1$.

\begin{lemma}[Local smoothness]\label[lemma]{lem:spikes:local}
Let $\Loss$ satisfy \cref{ass:spikes:L0L1}, let $\theta,u\in\R^d$, and put $g=\nabla\Loss(\theta)$ and $r=L_1\norm u$. Then
\begin{align}
  \norm{\nabla\Loss(\theta+u)-g}_*&\le\Bigl(\frac{L_0}{L_1}+\norm g_*\Bigr)\bigl(e^{r}-1\bigr),\label{eq:spikes:gradgrowth}\\
  \bigl\lvert\Loss(\theta+u)-\Loss(\theta)-\ip{g,u}\bigr\rvert&\le\Bigl(\frac{L_0}{L_1}+\norm g_*\Bigr)\frac{e^{r}-1-r}{L_1}.\label{eq:spikes:taylor}
\end{align}
Consequently, with $\varphi_1(c)=(e^c-1)/c$ and $\varphi_2(c)=(e^c-1-c)/c^2$, if $r\le c$ then
\begin{equation}\label{eq:spikes:local}
  \norm{\nabla\Loss(\theta+u)-g}_*\le\varphi_1(c)\,(L_0+L_1\norm g_*)\norm u,\qquad
  \Loss(\theta+u)\le\Loss(\theta)+\ip{g,u}+\varphi_2(c)\,(L_0+L_1\norm g_*)\norm u^2 .
\end{equation}
For $c=\frac13$, $\varphi_1(c)<1.19$ and $\varphi_2(c)<\frac35$. Both \eqref{eq:spikes:gradgrowth} and \eqref{eq:spikes:taylor} hold with equality for the cliff \eqref{eq:spikes:cliff} at $\theta\ge0$ and $u\ge0$.
\end{lemma}

As $c\to0$, $\varphi_1(c)\to1$ and $\varphi_2(c)\to\frac12$, and \eqref{eq:spikes:local} becomes the descent lemma for the local smoothness constant $L_0+L_1\norm g_*$ (\cref{lem:toolkit:descent}). The lemma says that this constant may be used as long as the step is short compared with $1/L_1$, at the price of the factors $\varphi_1,\varphi_2$.\sidenote{These are the functions $\varphi_1$ and $\varphi_2$ of exponential integrators for stiff differential equations, where they arise for the same reason: they integrate the exponential growth of a linear equation over one step.}

\begin{proof}
The plan is to bound the growth of the gradient along the segment from $\theta$ to $\theta+u$ by a linear differential inequality, solve it with Grönwall's inequality, and integrate. If $u=0$ there is nothing to prove, so let $u\ne0$. For $s\in[0,1]$ let $\phi(s)=\nabla\Loss(\theta+su)$, whose derivative $\phi'(s)=\nabla^2\Loss(\theta+su)u$ is continuous, and let $\psi(s)=\int_0^s\norm{\phi'(\sigma)}_*\dd\sigma$, a continuously differentiable function with $\psi(0)=0$. By the fundamental theorem of calculus and the triangle inequality for integrals, $\norm{\phi(s)-g}_*\le\psi(s)$. The assumption \eqref{eq:spikes:L0L1}, applied at $\theta+su$, bounds the derivative of $\psi$ by the current gradient, and the current gradient by $\norm g_*+\psi(s)$:
\[
  \psi'(s)=\norm{\nabla^2\Loss(\theta+su)u}_*\le\bigl(L_0+L_1\norm{\phi(s)}_*\bigr)\norm u\le a+r\,\psi(s),\qquad a=(L_0+L_1\norm g_*)\norm u .
\]
Hence $(e^{-rs}\psi(s))'=e^{-rs}(\psi'(s)-r\psi(s))\le ae^{-rs}$, and integrating from $0$ to $s$ gives $\psi(s)\le\frac ar(e^{rs}-1)$, the differential form of Grönwall's inequality (\cref{thm:toolkit:gronwall}). Since $a/r=L_0/L_1+\norm g_*$, the case $s=1$ is \eqref{eq:spikes:gradgrowth}. For \eqref{eq:spikes:taylor}, write the remainder of the first-order expansion as an integral and insert the bound on $\psi$:
\[
  \bigl\lvert\Loss(\theta+u)-\Loss(\theta)-\ip{g,u}\bigr\rvert=\Bigl\lvert\int_0^1\ip{\phi(s)-g,u}\dd s\Bigr\rvert\le\norm u\int_0^1\psi(s)\dd s\le\frac ar\norm u\Bigl(\frac{e^r-1}r-1\Bigr),
\]
which equals the right side of \eqref{eq:spikes:taylor}. The power series $\varphi_1(c)=\sum_{k\ge0}c^k/(k+1)!$ and $\varphi_2(c)=\sum_{k\ge0}c^k/(k+2)!$ have positive coefficients, so both functions increase on $(0,\infty)$; for $r\le c$ this gives $e^r-1=r\varphi_1(r)\le r\varphi_1(c)$ and $e^r-1-r=r^2\varphi_2(r)\le r^2\varphi_2(c)$, and substituting $r=L_1\norm u$ gives \eqref{eq:spikes:local}. The numbers: $\varphi_1(\frac13)=3(e^{1/3}-1)=1.187$, and $\varphi_2(\frac13)=9(e^{1/3}-\frac43)<9(\frac75-\frac43)=\frac35$ because $e<(\frac75)^3=2.744$. Finally, for the cliff at $\theta\ge0$ and $u\ge0$, both sides of \eqref{eq:spikes:gradgrowth} equal $\frac{L_0}{L_1}e^{L_1\theta}(e^{L_1u}-1)$, and integrating this in $u$ gives equality in \eqref{eq:spikes:taylor}.
\end{proof}

With the lemma, the main theorem is one application per step.

\begin{theorem}[Clipped steepest descent]\label{thm:spikes:clipped}
Let $\Loss$ satisfy \cref{ass:spikes:L0L1}, let $\theta_0\in\R^d$ and $\Delta_0=\Loss(\theta_0)-\Loss^\star$, and run \eqref{eq:spikes:clipped} with $0<\gamma\le1/(3L_1)$ and $0<\eta\le1/(2L_0)$ (every $\eta>0$ is allowed if $L_0=0$). Then:
\begin{enumerate}
\item $\Loss(\theta_{t+1})\le\Loss(\theta_t)-\frac12\min\bigl\{\eta\norm{g_t}_*^2,\ \gamma\norm{g_t}_*\bigr\}$ for every $t$;
\item for every $\eps>0$, the number of steps with $\norm{g_t}_*>\eps$ is at most $2\Delta_0\max\{1/(\eta\eps^2),1/(\gamma\eps)\}$. If $L_0>0$, then with $\eta=1/(2L_0)$ and $\gamma=1/(3L_1)$ it is at most
\begin{equation}\label{eq:spikes:clippedrate}
  \max\Bigl\{\frac{4L_0\Delta_0}{\eps^2},\ \frac{6L_1\Delta_0}{\eps}\Bigr\},
\end{equation}
which is $4L_0\Delta_0/\eps^2$ whenever $\eps\le2L_0/(3L_1)$.
\end{enumerate}
\end{theorem}

\begin{proof}
The plan is to apply \cref{lem:spikes:local} to one step, whose length is at most $\gamma$, and to check that the quadratic remainder costs at most half of the first-order decrease. If $g_t=0$, the step is zero and (i) holds with equality. Otherwise put $u=-\eta_tg_t^\sharp$, so that $\ip{g_t,u}=-\eta_t\norm{g_t}_*^2$ and $\norm u=\eta_t\norm{g_t}_*=\min\{\eta\norm{g_t}_*,\gamma\}$. Since $L_1\norm u\le L_1\gamma\le\frac13$, \eqref{eq:spikes:local} with $c=\frac13$ and $\varphi_2(\frac13)<\frac35$ gives
\[
  \Loss(\theta_{t+1})\le\Loss(\theta_t)-\eta_t\norm{g_t}_*^2+\tfrac35\bigl(L_0+L_1\norm{g_t}_*\bigr)\eta_t^2\norm{g_t}_*^2
  =\Loss(\theta_t)-\eta_t\norm{g_t}_*^2\Bigl(1-\tfrac35\bigl(L_0\eta_t+L_1\norm u\bigr)\Bigr).
\]
The two terms in the inner parenthesis are controlled by the two parameters: $L_0\eta_t\le L_0\eta\le\frac12$ and $L_1\norm u\le\frac13$. The large parenthesis is therefore at least $1-\frac35\cdot\frac56=\frac12$, and the decrease is at least $\frac12\eta_t\norm{g_t}_*^2=\frac12\min\{\eta\norm{g_t}_*^2,\gamma\norm{g_t}_*\}$, which is (i).

For (ii), sum (i) over $t<T$: the sum of the decreases is at most $\Loss(\theta_0)-\Loss(\theta_T)\le\Delta_0$. The functions $x\mapsto\eta x^2$ and $x\mapsto\gamma x$ increase, so every step with $\norm{g_t}_*>\eps$ decreases the loss by more than $\frac12\min\{\eta\eps^2,\gamma\eps\}$, and there are at most $2\Delta_0/\min\{\eta\eps^2,\gamma\eps\}=2\Delta_0\max\{1/(\eta\eps^2),1/(\gamma\eps)\}$ of them, for every $T$. Substituting $\eta=1/(2L_0)$ and $\gamma=1/(3L_1)$ gives \eqref{eq:spikes:clippedrate}; its first entry is the larger exactly when $\eps\le2L_0/(3L_1)$.
\end{proof}

The bound does not contain the largest gradient met along the way. A clipped step on a cliff is short, but the gradient there is large, and it decreases the loss by $\frac12\gamma\norm{g_t}_*$, more than an unclipped step does near the target; so the cliff costs nothing in \eqref{eq:spikes:clippedrate} beyond its share of $\Delta_0$. For gradient descent with a fixed step the largest gradient is the whole story. Let
\[
  M=\sup\bigl\{\norm{\nabla\Loss(\theta)}_*:\Loss(\theta)\le\Loss(\theta_0)\bigr\}
\]
be the largest gradient on the initial sublevel set, assumed finite.

\begin{corollary}[Gradient descent with a step set by $M$]\label[corollary]{cor:spikes:gd}
Let $\Loss$ satisfy \cref{ass:spikes:L0L1} with $M<\infty$, and let $0<\eta\le\min\{1/(2L_0),1/(3L_1M)\}$. Then unnormalized steepest descent $\theta_{t+1}=\theta_t-\eta g_t^\sharp$ keeps every iterate in the sublevel set, decreases the loss by at least $\frac\eta2\norm{g_t}_*^2$ at every step, and has $\norm{g_t}_*>\eps$ for at most $2\Delta_0/(\eta\eps^2)$ steps, for every $\eps\le M$. With the largest allowed $\eta$ this is $2\Delta_0\max\{2L_0,3L_1M\}/\eps^2$.
\end{corollary}

\begin{proof}
Put $\gamma=\eta M\le1/(3L_1)$ and argue by induction. If $\Loss(\theta_t)\le\Loss(\theta_0)$, then $\norm{g_t}_*\le M=\gamma/\eta$, so $\eta_t=\eta$ in \eqref{eq:spikes:clipped} and the step of gradient descent is the clipped step. \Cref{thm:spikes:clipped}(i) then gives a decrease of at least $\frac12\min\{\eta\norm{g_t}_*^2,\eta M\norm{g_t}_*\}=\frac\eta2\norm{g_t}_*^2$, so $\theta_{t+1}$ is in the sublevel set too. The count follows as in \cref{thm:spikes:clipped}(ii), since $\gamma\eps=\eta M\eps\ge\eta\eps^2$.
\end{proof}

So gradient descent works on the class if its step size is of order $1/(L_0+L_1M)$, and then it needs of order $(L_0+L_1M)\Delta_0/\eps^2$ steps. The next theorem shows that no fixed step size does much better: up to a logarithm, gradient descent must pay for $M$.

\begin{theorem}[A fixed step pays for the largest gradient]\label{thm:spikes:lower}
Let $L_0,L_1,M,\Delta_0>0$ and $0<\eps\le M/2$, with $\Delta_0\ge M/L_1$, and fix $x_0\in\R$. Let $\cF$ be the set of twice continuously differentiable functions $f\colon\R\to\R$ that are bounded below and $(L_0,L_1)$-smooth, with $f(x_0)-\inf f\le\Delta_0$ and $\abs{f'(x)}\le M$ whenever $f(x)\le f(x_0)$. For every step size $\eta>0$ there is an $f\in\cF$ on which gradient descent $x_{t+1}=x_t-\eta f'(x_t)$ has $\abs{f'(x_t)}>\eps$ for every $t<N$, with
\begin{equation}\label{eq:spikes:lower}
  N=\frac{L_1M}{\log(1+L_1M/L_0)}\cdot\frac{\Delta_0-4\eps^2/(\sqrt3L_0)}{8\eps^2}.
\end{equation}
If $\eta\ge\eta^*=2\log(1+L_1M/L_0)/(L_1M)$, then $f$ can be chosen so that $\abs{f'(x_t)}\ge M$ for every $t$.
\end{theorem}

When $\eps^2\le L_0\Delta_0/4$, the second factor of \eqref{eq:spikes:lower} is at least $\Delta_0/(20\eps^2)$. Because $\log(1+P)<P$ for $P>0$, the first factor exceeds $L_0$, so the theorem also contains the familiar lower bound of order $L_0\Delta_0/\eps^2$; what is new is the factor $L_1M/\log(1+L_1M/L_0)$, which is unbounded in $M$. The functions of the proof are functions of one variable; as functions of the first coordinate they give the same bound in $\R^d$ with the Euclidean norm.

\begin{proof}
The plan is a dichotomy: a large step diverges on a cliff, and a small step crawls on a plateau of slope just above $\eps$. Put $P=L_1M/L_0$ and $y_0=\log(1+P)$, so that $\eta^*=2y_0/(L_1M)$.

\emph{A large step diverges on the cliff.} Let $\eta\ge\eta^*$ and $f(x)=\spikescliff(x-c)$ with $c=x_0-y_0/L_1$. By \eqref{eq:spikes:cliff}, $f$ is $(L_0,L_1)$-smooth, $\abs{f'(x_0)}=\frac{L_0}{L_1}(e^{y_0}-1)=M$, the sublevel set $\{f\le f(x_0)\}$ is the interval $\abs{x-c}\le y_0/L_1$, on which $\abs{f'}\le M$, and $f(x_0)-\inf f=\frac{L_0}{L_1^2}(P-\log(1+P))<M/L_1\le\Delta_0$; so $f\in\cF$. In the variable $z_t=L_1(x_t-c)$ gradient descent reads $z_{t+1}=z_t\bigl(1-\eta L_0\varphi_1(\abs{z_t})\bigr)$, with the increasing function $\varphi_1$ of \cref{lem:spikes:local}. At $t=0$ we have $\eta L_0\varphi_1(y_0)=\eta L_0P/y_0=\eta L_1M/y_0\ge2$, hence $\abs{z_1}=\abs{z_0}(\eta L_0\varphi_1(\abs{z_0})-1)\ge\abs{z_0}$. Since $\varphi_1$ increases, the inequality $\eta L_0\varphi_1(\abs{z_t})\ge2$ propagates, and by induction $\abs{z_t}\ge y_0$ for every $t$, that is, $\abs{f'(x_t)}\ge M>\eps$.

\emph{A small step crawls on the plateau.} Let $\eta<\eta^*$. If $\Delta_0\le4\eps^2/(\sqrt3L_0)$, then $N\le0$ and there is nothing to prove, so assume the opposite, and let
\[
  \spikesplat(x)=\frac{4\eps^2}{L_0}\Bigl(\sqrt{1+\bigl(\tfrac{L_0x}{2\eps}\bigr)^2}-1\Bigr),\qquad
  \spikesplat'(x)=\frac{2\eps\,\sigma}{\sqrt{1+\sigma^2}},\qquad
  \spikesplat''(x)=\frac{L_0}{(1+\sigma^2)^{3/2}},\qquad \sigma=\frac{L_0x}{2\eps}.
\]
Then $\abs{\spikesplat'}<2\eps\le M$ and $0<\spikesplat''\le L_0$, so $f(x)=\spikesplat(x-c)$ is $(L_0,L_1)$-smooth for every $c$; choose $c<x_0$ with $f(x_0)=\Delta_0$, which puts $f$ in $\cF$. Solving $f(x_0)=\Delta_0$ with $a=L_0\Delta_0/(4\eps^2)$ gives $x_0-c=\frac{2\eps}{L_0}\sqrt{(1+a)^2-1}\ge\frac{2\eps}{L_0}a=\frac{\Delta_0}{2\eps}$. On the other hand $\abs{f'(x)}>\eps$ exactly when $\abs{x-c}>r_\eps=2\eps/(\sqrt3L_0)$. Each step moves the iterate by $\eta\abs{f'(x_t)}<2\eta\eps$, so $\abs{f'(x_t)}>\eps$ for every
\[
  t<\frac{x_0-c-r_\eps}{2\eta\eps}\ ,\qquad\text{and}\qquad
  \frac{x_0-c-r_\eps}{2\eta\eps}\ge\frac{\Delta_0-4\eps^2/(\sqrt3L_0)}{4\eta\eps^2}>\frac{\Delta_0-4\eps^2/(\sqrt3L_0)}{4\eta^*\eps^2}=N.\qedhere
\]
\end{proof}

\Textcite{ZhangHeSraJadbabaie2020} proved the first results of this kind: the lemma, with the factor $\varphi_1(c)$ in front of $L_1\norm g_*$ and a slightly larger one in front of $L_0$; a bound of the form $O(L_0\Delta_0/\eps^2+L_1^2\Delta_0/L_0)$ for clipped gradient descent, with larger absolute constants, under smoothness assumed on a neighbourhood of the sublevel set; and a lower bound of the form \eqref{eq:spikes:lower} for fixed steps, with $\log M+1$ in place of our logarithm. Our constants are smaller for two reasons: integrating the exact bound \eqref{eq:spikes:gradgrowth} gives the factor $\varphi_2(c)$, which is below $\frac12\varphi_1(c)$, in the descent inequality; and counting only the steps with $\norm{g_t}_*>\eps$ turns their sum of two terms into a maximum, which is never larger up to a constant, since $6L_1/\eps\le4L_0/\eps^2+\frac94L_1^2/L_0$ by the inequality of arithmetic and geometric means.

Normalized updates bound their step length by definition, so the same analysis applies to them, with the same weak dependence on $L_1$.

\begin{proposition}[Normalized steps are clipped by construction]\label[proposition]{prop:spikes:normalized}
Let $\Loss$ satisfy \cref{ass:spikes:L0L1} and $0<\gamma\le1/(3L_1)$.
\begin{enumerate}
\item Normalized steepest descent $\theta_{t+1}=\theta_t+\gamma\lmo(g_t)$ satisfies $\Loss(\theta_{t+1})\le\Loss(\theta_t)-\frac45\gamma\norm{g_t}_*+\frac35L_0\gamma^2$. Hence
\[
  \frac1T\sum_{t<T}\norm{\nabla\Loss(\theta_t)}_*\le\frac{5\Delta_0}{4\gamma T}+\frac{3L_0\gamma}4,
\]
and with $\gamma=\min\{1/(3L_1),2\eps/(3L_0)\}$ some $t<T$ has $\norm{\nabla\Loss(\theta_t)}_*\le\eps$ as soon as $T\ge\max\{15L_0\Delta_0/(4\eps^2),15L_1\Delta_0/(2\eps)\}$.
\item This holds for normalized gradient descent ($\ell_2$), for sign descent ($\ell_\infty$, with \eqref{eq:spikes:L0L1} read as $\norm{\nabla^2\Loss}_{\infty\to1}\le L_0+L_1\norm{\nabla\Loss}_1$), and for orthogonalized gradient descent on a weight matrix (the spectral norm, with the nuclear norm as dual).
\item Adam without bias correction, with $\epsilon=0$ and $\beta_1^2<\beta_2$, takes steps with $\norm{\theta_{t+1}-\theta_t}_\infty\le C\eta$, where $C=(1-\beta_1)/\sqrt{(1-\beta_2)(1-\beta_1^2/\beta_2)}$. If $\Loss$ satisfies \cref{ass:spikes:L0L1} in $\ell_\infty$ and $C\eta\le c/L_1$, every Adam step therefore satisfies \eqref{eq:spikes:local}.
\end{enumerate}
\end{proposition}

\begin{proof}
The plan is to repeat the proof of \cref{thm:spikes:clipped}(i) for a step of fixed length. For (i), $u=\gamma\lmo(g_t)$ has $\norm u\le\gamma$ and $\ip{g_t,u}=-\gamma\norm{g_t}_*$, also when $g_t=0$ and $\lmo(g_t)$ is an arbitrary point of the unit ball. \Cref{lem:spikes:local} with $c=\frac13$ gives
\[
  \Loss(\theta_{t+1})\le\Loss(\theta_t)-\gamma\norm{g_t}_*+\tfrac35(L_0+L_1\norm{g_t}_*)\gamma^2\le\Loss(\theta_t)-\tfrac45\gamma\norm{g_t}_*+\tfrac35L_0\gamma^2,
\]
using $\frac35L_1\gamma\le\frac15$. Summing over $t<T$ and dividing by $\frac45\gamma T$ gives the average. With the stated $\gamma$ the second term is at most $\eps/2$, and the first is at most $\eps/2$ once $T\ge5\Delta_0/(2\gamma\eps)$, which is the stated condition; the smallest of the $T$ gradients is at most their average. Part (ii) consists of instances of (i): the $\lmo$ of the three norms is $-g/\norm g_2$, $-\sign(g)$ and $-UV^\top$ (\cref{ch:steepest,ch:spectral}). For (iii), $\abs{m_{t,i}}/\sqrt{v_{t,i}}\le C$ for every coordinate by Chapter~2's bound on Adam's update (\cref{prop:sign:adambound}), so every step has $\ell_\infty$-length at most $C\eta$, and \cref{lem:spikes:local} applies to it.
\end{proof}

Part (iii) does not prove that Adam converges: its direction need not be a descent direction, and a rate needs the analyses of \cref{ch:adaptive}. It proves that Adam, like clipped descent, never takes a step that leaves the region where the local curvature bound holds, provided $C\eta\lesssim1/L_1$; for $(\beta_1,\beta_2)=(0.9,0.999)$ the constant is $C=7.27$, for $(0.9,0.95)$ it is $1.17$. \Cref{sec:spikes:anatomy} shows that Adam comes close to this bound exactly when a quiet block of weights wakes up.

The last result of the section is about noise. \Textcite{SimsekliSagunGurbuzbalaban2019} measured \term{heavy tails}\index{heavy-tailed noise} in the gradient noise of deep networks, and \textcite{ZhangEtAl2020Attention} found heavier tails for attention models than for convolutional networks and proposed clipping as the reason adaptive methods train them faster. If the noise has an infinite variance, stochastic gradient descent fails at the first step: on $\Loss(\theta)=\frac12\norm\theta_2^2$ with noise $\zeta_0$ of infinite second moment, $\E\Loss(\theta_1)=\frac12\E\norm{(1-\eta)\theta_0-\eta\zeta_0}_2^2=\infty$ for every $\eta>0$. Clipping restores a rate.

\begin{theorem}[Clipped SGD with heavy-tailed noise]\label{thm:spikes:heavy}
Let $\Loss$ be $L$-smooth in the Euclidean norm and bounded below, with $\Delta_0=\Loss(\theta_0)-\Loss^\star$, and let the stochastic gradients be unbiased, $\E[g_t\mid\theta_t]=\nabla\Loss(\theta_t)$, with a bounded moment of order $p\in(1,2]$: $\E[\norm{g_t}_2^p\mid\theta_t]\le\Gamma^p$. Clipped SGD $\theta_{t+1}=\theta_t-\eta\spikesclip_\tau(g_t)$ satisfies
\begin{equation}\label{eq:spikes:heavy}
  \frac1T\sum_{t<T}\E\norm{\nabla\Loss(\theta_t)}_2^2\le\frac{2\Delta_0}{\eta T}+\frac{\Gamma^{2p}}{\tau^{2(p-1)}}+L\eta\Gamma^p\tau^{2-p}.
\end{equation}
With $\tau=\Gamma T^{1/(3p-2)}$ and $\eta=(2\Delta_0/(L\Gamma^2))^{1/2}T^{-p/(3p-2)}$, the right side equals $(2\Gamma\sqrt{2L\Delta_0}+\Gamma^2)\,T^{-2(p-1)/(3p-2)}$, so $\min_{t<T}\E\norm{\nabla\Loss(\theta_t)}_2\lesssim T^{-(p-1)/(3p-2)}$.
\end{theorem}

\begin{proof}
The plan is to bound the bias and the second moment of a clipped gradient by the $p$th moment, and then to run the descent lemma. Fix $t$ and condition on $\theta_t$; write $g=g_t$ and $\hat g=\spikesclip_\tau(g)$. The bias: $\hat g$ and $g$ differ only where $\norm g_2>\tau$, by $\norm g_2-\tau<\norm g_2\le\norm g_2^p\tau^{1-p}$, so
\[
  \bigl\lVert\E\hat g-\nabla\Loss(\theta_t)\bigr\rVert_2\le\E\norm{\hat g-g}_2\le\frac{\Gamma^p}{\tau^{p-1}}\eqqcolon b .
\]
The second moment: $\norm{\hat g}_2^2=\min\{\norm g_2,\tau\}^2\le\norm g_2^p\tau^{2-p}$, by checking the cases $\norm g_2\le\tau$ and $\norm g_2>\tau$, so $\E\norm{\hat g}_2^2\le\Gamma^p\tau^{2-p}$. The descent lemma (\cref{lem:toolkit:descent}) gives $\Loss(\theta_{t+1})\le\Loss(\theta_t)-\eta\ip{\nabla\Loss(\theta_t),\hat g}+\frac L2\eta^2\norm{\hat g}_2^2$. The bias enters only the middle term, through $\ip{\nabla\Loss,\E\hat g}\ge\norm{\nabla\Loss}_2^2-\norm{\nabla\Loss}_2b\ge\frac12\norm{\nabla\Loss}_2^2-\frac12b^2$. Taking expectations,
\[
  \E\Loss(\theta_{t+1})\le\E\Loss(\theta_t)-\frac\eta2\E\norm{\nabla\Loss(\theta_t)}_2^2+\frac\eta2b^2+\frac L2\eta^2\Gamma^p\tau^{2-p}.
\]
Summing over $t<T$, using $\E\Loss(\theta_T)\ge\Loss^\star$ and dividing by $\eta T/2$ gives \eqref{eq:spikes:heavy}. With the stated $\tau$ and $\eta$, the first and third terms both equal $\Gamma\sqrt{2L\Delta_0}\,T^{-2(p-1)/(3p-2)}$ and the second equals $\Gamma^2T^{-2(p-1)/(3p-2)}$. The last claim follows because the smallest of the $T$ numbers $(\E\norm{\nabla\Loss(\theta_t)}_2)^2\le\E\norm{\nabla\Loss(\theta_t)}_2^2$ is at most their average.
\end{proof}

For $p=2$ the rate is $T^{-1/4}$, the rate of stochastic gradient descent with bounded variance; for $p\to1$ it degrades to nothing. The threshold must grow with the length of the run, like $T^{1/(3p-2)}$: a fixed threshold leaves the bias $\Gamma^p/\tau^{p-1}$, and clipping with noise is biased (\cref{ex:spikes:bias}). \Textcite{ZhangEtAl2020Attention} proved this rate and a matching lower bound, and \textcite{KoloskovaHendrikxStich2023} analysed the bias of a fixed threshold.

\section{The theory: the anatomy of a spike}\label{sec:spikes:anatomy}

A \term{loss spike}\index{loss spike} is a step too long for the curvature it lands in. \Cref{thm:spikes:clipped} says that a step of bounded length is safe; so the question is how a clipped, well-tuned run produces steps that are too long, or curvature that is too high. We take apart three mechanisms, one for each further defence in the practice box: Adam's state ($\beta_2$), the output logits (z-loss) and the attention logits (QK-norm, MuonClip).

\paragraph{The optimizer's state.} \Cref{ch:adaptive} showed that after a long stretch of tiny gradients a coordinate whose gradient jumps receives an Adam step of about $(1-\beta_1)/\sqrt{1-\beta_2}$ times the learning rate (\cref{prop:adaptive:quietloud}). That is the first step of a burst, not its peak.\sidenote{\Textcite{KingmaBa2015} observed that Adam's step is roughly bounded by the learning rate, and singled out the most severe case of sparsity, a gradient that was zero at all earlier steps, in which the step is $(1-\beta_1)/\sqrt{1-\beta_2}$ times the learning rate.}

\begin{proposition}[A quiet block wakes up]\label[proposition]{prop:spikes:burst}
Consider one coordinate of Adam with $\beta_1,\beta_2\in[0,1)$ and $\epsilon\ge0$, moments $m_t=\beta_1m_{t-1}+(1-\beta_1)g_t$ and $v_t=\beta_2v_{t-1}+(1-\beta_2)g_t^2$, and step $-\eta u_t$ with $u_t=m_t/(\sqrt{v_t}+\epsilon)$. Suppose that $\abs{m_{t_0-1}}\le\delta$ and $0\le v_{t_0-1}\le\delta^2$, as holds when $\abs{g_s}\le\delta$ for every $s<t_0$ and $m_{-1}=v_{-1}=0$, and that $g_s=\bar g\ne0$ for $t_0\le s<t_0+k$. At the $k$th loud step $t=t_0+k-1$,
\begin{equation}\label{eq:spikes:burstbounds}
  \sign(\bar g)\,u_t\le\frac{(1-\beta_1^k)\abs{\bar g}+\beta_1^k\delta}{\sqrt{1-\beta_2^k}\,\abs{\bar g}+\epsilon},
\end{equation}
and $\sign(\bar g)u_t\ge\bigl((1-\beta_1^k)\abs{\bar g}-\beta_1^k\delta\bigr)/\bigl(\sqrt{1-\beta_2^k}\,\abs{\bar g}+\beta_2^{k/2}\delta+\epsilon\bigr)$ whenever the numerator is nonnegative. If the coordinate was silent, $\delta=0$, then $u_t=\sign(\bar g)\,a_k$ with
\begin{equation}\label{eq:spikes:burst}
  a_k=\frac{1-\beta_1^k}{\sqrt{1-\beta_2^k}+\epsilon/\abs{\bar g}} .
\end{equation}
With $\epsilon=0$: $a_1=(1-\beta_1)/\sqrt{1-\beta_2}$; $a_k\to1$ as $k\to\infty$; $a_k\le C(\beta_1,\beta_2)=(1-\beta_1)/\sqrt{(1-\beta_2)(1-\beta_1^2/\beta_2)}$ for every $k$ when $\beta_1^2<\beta_2$; and $a_k\ge(1-\beta_1^k)/\sqrt{k(1-\beta_2)}$, so that for $\beta_1=0.9$, $\max_ka_k\ge a_{12}\ge0.207/\sqrt{1-\beta_2}$.
\end{proposition}

\begin{proof}
The plan is to unroll the two recursions over the loud steps. Over $k$ steps with the constant gradient $\bar g$, the first moment forgets its initial value at the rate $\beta_1$ and the second at the rate $\beta_2$:
\[
  m_t=\beta_1^km_{t_0-1}+(1-\beta_1^k)\bar g,\qquad v_t=\beta_2^kv_{t_0-1}+(1-\beta_2^k)\bar g^2,
\]
by induction on $k$, using $(1-\beta)\sum_{j<k}\beta^j=1-\beta^k$. Hence $\abs{\sign(\bar g)m_t-(1-\beta_1^k)\abs{\bar g}}\le\beta_1^k\delta$ and $(1-\beta_2^k)\bar g^2\le v_t\le(1-\beta_2^k)\bar g^2+\beta_2^k\delta^2$. For the upper bound, if $\sign(\bar g)m_t<0$ there is nothing to prove; otherwise divide the upper bound on $\sign(\bar g)m_t$ by the lower bound $\sqrt{1-\beta_2^k}\abs{\bar g}+\epsilon$ on $\sqrt{v_t}+\epsilon$. For the lower bound, divide the nonnegative lower bound on $\sign(\bar g)m_t$ by the upper bound on $\sqrt{v_t}+\epsilon$, using $\sqrt{x+y}\le\sqrt x+\sqrt y$. With $\delta=0$ the two bounds coincide and give \eqref{eq:spikes:burst}. The limit $a_k\to1$ is immediate. The bound by $C(\beta_1,\beta_2)$ is Chapter~2's bound on $\abs{m_t}/\sqrt{v_t}$ (\cref{prop:sign:adambound}), which holds for every sequence of gradients, here the sequence that is zero up to $t_0-1$ and $\bar g$ afterwards. Finally $1-\beta_2^k=(1-\beta_2)\sum_{j<k}\beta_2^j\le k(1-\beta_2)$, and $(1-0.9^{12})/\sqrt{12}=0.2071$.
\end{proof}

In words: after a long silence both moments are empty, and the coordinate goes through a \term{burst}\index{Adam!burst after a quiet period}. The first moment refills in about $1/(1-\beta_1)=10$ steps, the second in about $1/(1-\beta_2)$ steps, and in between $m$ has nearly its full size while $\sqrt v$ does not. The step grows for about $1.26/\log(1/\beta_1)\approx12$ steps (the maximizer of $(1-\beta_1^k)/\sqrt k$), and then decays to $1$ over about $1/(1-\beta_2)$ steps, as the second moment catches up:

\begin{center}\small
\begin{tabular}{@{}lccc@{}}
\toprule
$\beta_2$ & first step $a_1$ & peak $\max_ka_k$ (at $k$) & worst case $C(0.9,\beta_2)$\\
\midrule
$0.95$ & $0.45$ & $1.10$ ($k=20$) & $1.17$\\
$0.99$ & $1.00$ & $2.13$ ($k=13$) & $2.35$\\
$0.999$ & $3.16$ & $6.57$ ($k=12$) & $7.27$\\
\bottomrule
\end{tabular}
\end{center}

\noindent The burst reaches about nine tenths of the worst case over all gradient sequences, which is approached by gradients that grow geometrically. Bias correction does not help: it divides by $1-\beta_1^{t+1}$ and $1-\beta_2^{t+1}$ with $t$ counted from the start of training, and for $t\gg1/(1-\beta_2)$ both are $1$. The burst is precisely the transient that bias correction removes at the start of training (\cref{ch:warmup}), recreated for a coordinate whose memory has been emptied at a time when bias correction has long stopped acting; a bias correction that counted the age $k$ of the coordinate would remove it (\cref{ex:spikes:burst}). Without momentum the burst is a single step: for RMSProp, $\beta_1=0$, the largest $a_k$ is $a_1=1/\sqrt{1-\beta_2}$, which is $31.6$ at $\beta_2=0.999$. And $\epsilon$ caps the burst only for gradients with $\abs{\bar g}\lesssim\epsilon/\sqrt{1-\beta_2^k}$.

Which blocks go silent? Rows of an embedding matrix for tokens absent from recent batches, experts that a router stops choosing, units whose inputs vanish; in a milder form, every block whose gradient grows by a large factor within fewer than $1/(1-\beta_2)$ steps (\cref{ex:spikes:jump}). A whole block that wakes up together moves, coordinate by coordinate, by up to $6.6$ learning rates per step at $\beta_2=0.999$, against about one for steady gradients; whether the loss then spikes depends on the curvature of the block, which the proposition does not see.

\paragraph{The output logits.} The second mechanism is a direction in which the loss does not change at all.

\begin{proposition}[A flat direction and the z-loss]\label[proposition]{prop:spikes:zloss}
For logits $z\in\R^K$ let $Z(z)=\sum_ke^{z_k}$ and $\pi=\softmax(z)$, and let $\ell(z)=-\sum_ky_kz_k+\log Z(z)$ be the cross-entropy against a probability vector $y$ (a one-hot label, or the label frequencies of a batch).
\begin{enumerate}
\item $\ell(z+c\mathbf1)=\ell(z)$ for every $c\in\R$; $\ip{\nabla\ell(z),\mathbf1}=0$ and $\nabla^2\ell(z)\mathbf1=0$.
\item For an output layer $z=Wh+b$ with $W\in\R^{K\times d_{\mathrm{in}}}$, every gradient of the loss of a batch satisfies $\mathbf1^\top\nabla_W\ell=0$ and $\ip{\mathbf1,\nabla_b\ell}=0$. Hence an update that is a linear combination of gradients, multiplied on the right by matrices (gradient descent, momentum, right preconditioning), and also an orthogonalized update $-\eta\polar(G)$, leaves $\mathbf1^\top W$ and $\ip{\mathbf1,b}$ unchanged; a sign step on $b$ changes $\ip{\mathbf1,b}$ by $-\eta\sum_k\sign(\pi_k-y_k)$.
\item With the z-loss, $\ell_\alpha(z)=\ell(z)+\alpha(\log Z(z))^2$ with $\alpha>0$, the restriction to the flat line is a parabola: $\ell_\alpha(z+c\mathbf1)=\ell(z)+\alpha(\log Z(z)+c)^2$, with second derivative $2\alpha$ in $c$ and minimum at $\log Z=0$. If $y_k>0$ for every $k$, the only stationary point of $\ell_\alpha$ is $z=\log y$ (coordinatewise), where $\log Z=0$ and $\pi=y$.
\end{enumerate}
\end{proposition}

\begin{proof}
The plan is to compute the gradient and the Hessian of $\ell$ and to read everything off them. (i) $Z(z+c\mathbf1)=e^cZ(z)$, so $\log Z$ increases by $c$, exactly as much as $\sum_ky_kz_k$ does. Differentiating, $\nabla\ell=\pi-y$, whose coordinates sum to $1-1=0$, and $\nabla^2\ell=\diag(\pi)-\pi\pi^\top$, which maps $\mathbf1$ to $\pi-\pi=0$. (ii) For one example the chain rule gives $\nabla_W\ell=(\pi-y)h^\top$ and $\nabla_b\ell=\pi-y$, so $\mathbf1^\top\nabla_W\ell=0$ by (i); a batch gradient is an average of such terms. A linear combination $\sum_sc_sG_sR_s$ inherits $\mathbf1^\top G_s=0$. If $\mathbf1^\top G=0$, then $\mathbf1$ is orthogonal to the column space of $G$, which is spanned by the left singular vectors $U$, so $\mathbf1^\top\polar(G)=\mathbf1^\top UV^\top=0$. The sign step is a computation. (iii) The first claim follows from $\log Z(z+c\mathbf1)=\log Z(z)+c$ and (i). A stationary point satisfies $\nabla\ell_\alpha=\pi-y+2\alpha\log Z\,\pi=0$, that is, $(1+2\alpha\log Z)\pi=y$; summing the coordinates gives $1+2\alpha\log Z=1$, hence $\log Z=0$ and $\pi=y$, and then $z_k=\log\pi_k+\log Z=\log y_k$.
\end{proof}

Gradient descent therefore never moves the logits along this \term{flat direction}\index{flat direction of the softmax}, and neither does Muon (\cref{ex:spikes:muonmean}), whereas coordinatewise normalization does:\sidenote{In the recipes of \textcite{JordanEtAl2024Muon} and \textcite{LiuEtAl2025Muon}, Muon trains the hidden matrices and AdamW trains the embedding and output layers, so the output logits are still exposed to the drift.} the coordinates of $\pi-y$ sum to zero, their signs need not. Decoupled weight decay shrinks the flat component, and the z-loss adds a spring of stiffness $2\alpha$ that holds $\log Z$ near $0$. In exact arithmetic the drift is harmless, since the loss does not see it; it matters through finite precision, for instance in bfloat16, which keeps $8$ significant bits, so that logits near $100$ are rounded to multiples of $0.5$ and the probabilities computed from them are off by factors up to $e^{\pm0.25}$.\sidenote{\Textcite{ZophEtAl2022} introduced a z-loss on the logits of the router of a mixture of experts, motivated by roundoff in the exponentials of large logits. Whether the output-logit divergence of \cref{obs:spikes:proxies} is a precision effect has, as far as we know, not been established.} How strong a spring is needed depends on how fast the drift is, which the proposition does not say; \cref{sec:spikes:reading} measures it in a small model.

\paragraph{The attention logits.} The third mechanism is a scale, that of the \term{attention logits}\index{attention logits}.

\begin{proposition}[Attention logits, QK-norm and QK-clip]\label[proposition]{prop:spikes:qknorm}
For an attention head with query and key weights $W_Q,W_K\in\R^{d_{\mathrm h}\times d_{\mathrm{in}}}$ and inputs $x,x'\in\R^{d_{\mathrm{in}}}$, let $s(x,x')=\ip{W_Qx,W_Kx'}/\sqrt{d_{\mathrm h}}$.
\begin{enumerate}
\item The logit is bilinear in $(W_Q,W_K)$: replacing $(W_Q,W_K)$ by $(\rho W_Q,\rho W_K)$ multiplies every logit by $\rho^2$, and $\abs{s(x,x')}\le\norm{W_Q}_{\op}\norm{W_K}_{\op}\norm x_2\norm{x'}_2/\sqrt{d_{\mathrm h}}$.
\item With QK-norm, $\tilde s(x,x')=\ip{\omega_q\odot\mathcal N(W_Qx),\,\omega_k\odot\mathcal N(W_Kx')}/\sqrt{d_{\mathrm h}}$, where $\mathcal N$ is RMS normalization or a LayerNorm without bias and without $\epsilon$, and $\omega_q,\omega_k\in\R^{d_{\mathrm h}}$ are gains,
\[
  \abs{\tilde s(x,x')}\le\sqrt{d_{\mathrm h}}\,\norm{\omega_q}_\infty\norm{\omega_k}_\infty
\]
for all $W_Q,W_K,x,x'$ for which the normalizations are defined, with equality when $\mathcal N(W_Qx)=\mathcal N(W_Kx')$ and both gains are the same constant vector.
\item If the logits of one row of attention lie in an interval of length $R$, then each of the $n_{\mathrm{key}}$ attention weights of the row is at most $1/(1+(n_{\mathrm{key}}-1)e^{-R})$, and the row's entropy is at least $\log(1+(n_{\mathrm{key}}-1)e^{-R})$.
\item (QK-clip.) If the largest logit of a head on a batch is $s_{\max}>\bar s>0$, then replacing $(W_Q,W_K)$ by $(\rho W_Q,\rho W_K)$ with $\rho=\sqrt{\bar s/s_{\max}}$ multiplies every logit of the head by $\bar s/s_{\max}$, so that the largest logit on that batch becomes exactly $\bar s$.
\end{enumerate}
\end{proposition}

\begin{proof}
The plan is to treat the four parts in turn; only (ii) and (iii) need a computation. (i) Bilinearity is immediate, and Cauchy--Schwarz gives $\abs{\ip{W_Qx,W_Kx'}}\le\norm{W_Qx}_2\norm{W_Kx'}_2\le\norm{W_Q}_{\op}\norm x_2\norm{W_K}_{\op}\norm{x'}_2$. (ii) RMS normalization maps $v$ to $v/(\norm v_2/\sqrt{d_{\mathrm h}})$, and a LayerNorm without bias and $\epsilon$ maps it to $(v-\bar v\mathbf1)/(\norm{v-\bar v\mathbf1}_2/\sqrt{d_{\mathrm h}})$, with $\bar v$ the mean of the coordinates; both images have Euclidean norm $\sqrt{d_{\mathrm h}}$. Hence, with $a=\mathcal N(W_Qx)$ and $b=\mathcal N(W_Kx')$,
\[
  \Bigl\lvert\sum_i\omega_{q,i}\omega_{k,i}a_ib_i\Bigr\rvert\le\norm{\omega_q}_\infty\norm{\omega_k}_\infty\sum_i\abs{a_ib_i}\le\norm{\omega_q}_\infty\norm{\omega_k}_\infty\norm a_2\norm b_2=\norm{\omega_q}_\infty\norm{\omega_k}_\infty\,d_{\mathrm h},
\]
and dividing by $\sqrt{d_{\mathrm h}}$ gives the bound; in the stated case every inequality is an equality. (iii) Let the logits be $s_1,\dots,s_{n_{\mathrm{key}}}$. Every $s_i$ is at least $s_j-R$, so the weight of key $j$ is $e^{s_j}/\sum_ie^{s_i}\le e^{s_j}/(e^{s_j}+(n_{\mathrm{key}}-1)e^{s_j-R})$, which is the bound; the entropy $\sum_jp_j\log(1/p_j)$ is at least $\log(1/\max_jp_j)$. (iv) follows from (i).
\end{proof}

Part (i) is the source of the growth: the logit is quadratic in the scale of $(W_Q,W_K)$, and under Adam the scale can grow linearly in time. Every coordinate moves by at most $C\eta$ per step (\cref{prop:spikes:normalized}(iii)), so after $T$ steps $\norm{W_Q-W_Q^{(0)}}_{\op}\le\norm{W_Q-W_Q^{(0)}}_F\le C\eta T\sqrt{d_{\mathrm h}d_{\mathrm{in}}}$, and for inputs normalized to $\norm x_2=\sqrt{d_{\mathrm{in}}}$ the bound of (i) grows like $(C\eta T)^2d_{\mathrm{in}}^2\sqrt{d_{\mathrm h}}$. The Frobenius bound is attained by rank-one sign updates $-\eta\sign(\delta)\sign(x)^\top$, which is what Adam does with a consistent rank-one gradient. Under QK-norm the same argument applied to the gains gives $\sqrt{d_{\mathrm h}}(\norm{\omega^{(0)}}_\infty+C\eta T)^2$: the growth rate of the worst case falls by the factor $d_{\mathrm{in}}^2$. QK-norm does not forbid sharp attention; with unit gains and $d_{\mathrm h}=128$ the logits may range over an interval of length $2\sqrt{128}\approx23$, enough for nearly one-hot rows. What it removes is growth driven by the scale of the weight matrices, and part (iii) says that sharp attention needs large logits: entropy collapse in the sense of \textcite{ZhaiEtAl2023} is a growth of the logit range.\sidenote{\Textcite{HenryEtAl2020} normalized queries and keys to unit Euclidean norm with a learned temperature for machine translation; \textcite{DehghaniEtAl2023} used LayerNorm on queries and keys to remove the divergence of ViT-22B.} QK-clip acts on the same scale after the fact: by (iv) it is the smallest rescaling of $(W_Q,W_K)$, split evenly between them, that brings the largest logit of the batch back to the threshold.

\begin{fieldguide}{MuonClip}{\textcite{KimiTeam2025}}
\fgrow{update}{a Muon step on every hidden matrix (\cref{ch:spectral}); then, for each attention head whose largest logit $s_{\max}$ on the batch exceeds $\bar s$, $W_Q\leftarrow\rho W_Q$ and $W_K\leftarrow\rho W_K$ with $\rho=\sqrt{\bar s/s_{\max}}$}
\fgrow{norm}{spectral, for the Muon step}
\fgrow{preconditioner}{none beyond orthogonalization}
\fgrow{averages}{Muon's momentum}
\fgrow{state}{one momentum matrix per weight; the largest logit of each head, read off the forward pass}
\fgrow{in the wild}{$\bar s=100$ in Kimi K2; Muon itself as scaled by \textcite{LiuEtAl2025Muon}}
\fgrow{caution}{acts on the weights after the step, not on the gradient, and only on heads above the threshold}
\end{fieldguide}

\section{Reading the theory}\label{sec:spikes:reading}

\paragraph{The price of a fixed step.} \Cref{fig:spikes:separation} puts the two theorems side by side on the simplest function that has both a cliff and a plateau, $f(x,y)=\spikescliff(x)+\spikesplat(y)$ with $L_0=L_1=1$, which is $(1,1)$-smooth in the Euclidean norm because its Hessian is diagonal. The largest step with which gradient descent does not diverge is exactly $\eta^*=2\log(1+L_1M/L_0)/(L_1M)$, the cliff's threshold from the proof of \cref{thm:spikes:lower}: $1.39$ at $M=1$, $0.092$ at $M=100$ and $2.8\cdot10^{-5}$ at $M=10^6$. Gradient descent at $0.9\eta^*$ then crawls along the plateau, and its number of steps equals the gradient-flow time of the plateau divided by the step size to within one step, because the steps are small compared with $1/L_0$: from $81$ steps at $M=1$ to $4.07\cdot10^6$ at $M=10^6$, growing like $M/\log M$. Clipped descent with the parameters of \cref{thm:spikes:clipped} needs $203$ steps at $M=1$ and $241$ at $M=10^6$: the plateau costs it $2\cdot101$ steps at the step size $1/(2L_0)$, and the cliff $3\log(1+M)$ normalized steps of length $\gamma=1/3$. Its largest convergent step size does not depend on $M$ at all. For every start it lies between $L_1\gamma/(L_0(e^{L_1\gamma/2}-1))=1.838$, below which clipped descent cannot get caught in a two-cycle across the bottom of the valley, and $2/L_0$, above which it cannot converge at the bottom; the measured values fill this band.

\begin{figure}[t]
\begin{fullwidth}
\centering
\begin{tikzpicture}
\begin{groupplot}[group style={group size=2 by 1,horizontal sep=18mm},wplot,width=82mm,height=52mm,xmode=log,xmin=0.7,xmax=1.5e6]
\nextgroupplot[ymode=log,xlabel={largest gradient $M$},ylabel={steps to $\norm{\nabla f}_2\le\eps$},ymin=40,ymax=1e7,legend style={at={(0.03,0.97)},anchor=north west,font=\scriptsize}]
\addplot[theory] table[x=M,y=gdtheory] {figures/data/10-spikes-separation.dat};
\addplot[simulation] table[x=M,y=gdsim] {figures/data/10-spikes-separation.dat};
\addplot[theory2] table[x=M,y=cliptheory] {figures/data/10-spikes-separation.dat};
\addplot[simulation2] table[x=M,y=clipsim] {figures/data/10-spikes-separation.dat};
\legend{gradient flow$/\eta$,{GD, $\eta=0.9\eta^*$},flow$/\eta$ $+$ cliff steps,clipped GD}
\nextgroupplot[ymode=log,xlabel={largest gradient $M$},ylabel={largest stable step size},ymin=1e-5,ymax=5,legend style={at={(0.03,0.03)},anchor=south west,font=\scriptsize}]
\fill[seamist] (axis cs:0.7,1.83796) rectangle (axis cs:1.5e6,2);
\addplot[theory] table[x=M,y=thrtheory] {figures/data/10-spikes-separation.dat};
\addplot[simulation] table[x=M,y=thrsim] {figures/data/10-spikes-separation.dat};
\addplot[simulation2] table[x=M,y=clipthr] {figures/data/10-spikes-separation.dat};
\legend{$\eta^*=2\log(1+M)/M$,GD (bisection),clipped GD (scan)}
\end{groupplot}
\end{tikzpicture}
\end{fullwidth}
\caption{Clipped and fixed-step gradient descent on a cliff followed by a plateau, $f(x,y)=\spikescliff(x)+\spikesplat(y)$ with $L_0=L_1=1$ and $\eps=0.05$, started where the slope of the cliff is $M$ and the height of the plateau is $1$. Left: steps until the gradient norm is at most $\eps$, for gradient descent at $0.9$ times its largest stable step (circles) and clipped descent with $\eta=1/(2L_0)$, $\gamma=1/(3L_1)$ (squares); curves: the gradient-flow time of the plateau, $101.15$, divided by the step size, plus for clipped descent the $3\log(1+M)-3\log\frac53$ normalized steps down the cliff. Right: the largest step size with which gradient descent does not diverge (circles, bisection) against $\eta^*$ (curve), and the largest with which clipped descent converges (squares, scan in steps of $0.0025$), against the band $[1.838,2]$ predicted for every start (shaded). Deterministic; script \texttt{10-spikes-separation.py}.}
\label{fig:spikes:separation}
\end{figure}

The agreement is exact because the functions are those of the proofs; it is a check of the arithmetic, not of the idealization. It fails as soon as the step sizes stop being small compared with $1/L_0$ on the plateau, where the flow approximation of the step count breaks down, and it says nothing about functions that are not separable.

\paragraph{Where to clip.} In units of the gradient norm, the threshold of \cref{thm:spikes:clipped} is $\gamma/\eta=2L_0/(3L_1)$: clipping should begin where the term $L_1\norm{\nabla\Loss}_*$ of the curvature bound starts to dominate $L_0$. That is a property of the loss, not of the optimizer, and the theory gives no reason for a universal value such as $1.0$; global gradient norms change with the normalization of the loss and with the number of parameters. What it predicts is that a good threshold is inactive where the landscape is flat and active on cliffs and during spikes, and that the price of a threshold set too low is the factor $1/(\gamma\eps)$ of \cref{thm:spikes:clipped}(ii): the run is slowed wherever the gradient exceeds the threshold.

\paragraph{Clipping in front of Adam.} Adam's update is unchanged when all gradients are multiplied by the same factor (up to $\epsilon$), and \cref{prop:spikes:normalized}(iii) bounds its step whatever the gradient. So clipping in front of Adam does not bound Adam's step; it bounds the weight that one large gradient receives in $m$ and $v$. When the global gradient norm jumps by a large factor at one step and the threshold sits near its old value, the clipped gradient keeps the old norm, and the loud event that would start a burst is suppressed. A block that wakes up while the global norm stays below the threshold is not affected: the mechanism of \cref{prop:spikes:burst} survives global clipping.

\begin{figure}[t]
\begin{fullwidth}
\centering
\begin{tikzpicture}
\begin{groupplot}[group style={group size=2 by 1,horizontal sep=18mm},wplot,width=82mm,height=52mm]
\nextgroupplot[xmode=log,xmin=1,xmax=3000,ymin=0,ymax=7.8,xlabel={steps $k$ since the block woke up},ylabel={Adam step $\abs{\Delta\theta_i}/\eta$},legend style={at={(0.97,0.97)},anchor=north east,font=\scriptsize},legend columns=2]
\addplot[guide,domain=1:3000] {7.2703};
\addplot[guide,domain=1:3000] {2.3452};
\addplot[guide,domain=1:3000] {1.1650};
\addplot[theory3] table[x=k,y=th999] {figures/data/10-spikes-burst-k.dat};
\addplot[simulation,each nth point=2] table[x=k,y=sim999] {figures/data/10-spikes-burst-k.dat};
\addplot[theory2] table[x=k,y=th99] {figures/data/10-spikes-burst-k.dat};
\addplot[simulation2,each nth point=2] table[x=k,y=sim99] {figures/data/10-spikes-burst-k.dat};
\addplot[theory] table[x=k,y=th95] {figures/data/10-spikes-burst-k.dat};
\addplot[simulation,mark=triangle*,mark size=1.8pt,each nth point=2] table[x=k,y=sim95] {figures/data/10-spikes-burst-k.dat};
\legend{,,,{$a_k$, $\beta_2=0.999$},Adam,{$a_k$, $\beta_2=0.99$},Adam,{$a_k$, $\beta_2=0.95$},Adam}
\nextgroupplot[xmode=log,ymode=log,x dir=reverse,xmin=7e-5,xmax=0.14,ymin=0.25,ymax=30,xlabel={$1-\beta_2$},ylabel={largest step $/\eta$},legend style={at={(0.03,0.97)},anchor=north west,font=\scriptsize}]
\addplot[guide] table[x=omb2,y=worst] {figures/data/10-spikes-burst-peak.dat};
\addplot[theory] table[x=omb2,y=peak] {figures/data/10-spikes-burst-peak.dat};
\addplot[theory3,densely dotted] table[x=omb2,y=a1] {figures/data/10-spikes-burst-peak.dat};
\addplot[theory2] table[x=omb2,y=peakeps] {figures/data/10-spikes-burst-peak.dat};
\addplot[simulation] table[x=omb2,y=sim] {figures/data/10-spikes-burst-peak.dat};
\addplot[simulation2] table[x=omb2,y=simeps] {figures/data/10-spikes-burst-peak.dat};
\legend{{worst case $C(\beta_1,\beta_2)$},peak $\max_ka_k$,first step $a_1$,{peak, $\abs{\bar g}=10\epsilon$},Adam,{Adam, $\abs{\bar g}=10\epsilon$}}
\end{groupplot}
\end{tikzpicture}
\end{fullwidth}
\caption{Adam's step on a block of $32$ weights of a linear regression whose input features are absent for the first $10^5$ steps of training, so that the block's moments are zero, and present afterwards ($\beta_1=0.9$, $\epsilon=10^{-8}$, bias correction counted from the start of training, batches of $4096$, step size $10^{-6}$, labels $\ip{\theta_\star,x}$ plus unit Gaussian noise with $\theta_\star\in\{\pm1\}^{32}$). Left: the median over the block of $\abs{\Delta\theta_i}/\eta$ at the $k$th step after the features appear (points, every other measured $k$ shown), against $a_k$ of \eqref{eq:spikes:burst} (curves); dashed lines: the worst case $C(0.9,\beta_2)$. Right: the largest such step over the first $80$ steps against $1-\beta_2$ (points), with the burst peak $\max_ka_k$, the first step $a_1$ and the worst case; the lower curve and squares repeat the computation with the loss multiplied by $10^{-7}$, so that $\abs{\bar g}=10\epsilon$. One run per point; seeds in \texttt{10-spikes-burst.py}.}
\label{fig:spikes:burst}
\end{figure}

\paragraph{The burst.} \Cref{fig:spikes:burst} shows the anatomy of \cref{prop:spikes:burst} in a training run. The measured steps follow $a_k$ to within one percent: they rise for about twelve steps, peak at $6.59$, $2.13$ and $1.09$ learning rates for $\beta_2=0.999,0.99,0.95$ (predicted $6.57$, $2.13$, $1.10$), and decay to $1$ over about $1/(1-\beta_2)$ steps; the small differences come from minibatch noise, which enters $v$ but not, on average, $m$. Across $\beta_2$ the peak follows $0.207/\sqrt{1-\beta_2}$ for long memories, stays below the worst case, and is roughly twice the single-jump value $a_1$ of \cref{ch:adaptive}. Lowering $\beta_2$ from $0.999$ to $0.95$ shrinks the largest burst six-fold, from $6.6$ to $1.1$ learning rates per coordinate. This is a mechanism behind the choice $\beta_2=0.95$ in large language models and behind \cref{obs:spikes:stale}; that it is \emph{the} reason for the choice is a suggestion, not a result, and the price of a short memory is a noisier preconditioner (\cref{ch:adaptive}). A larger $\epsilon$ also damps the burst, but only for coordinates whose gradient is not much larger than $\epsilon/\sqrt{1-\beta_2^k}$: at $\abs{\bar g}=10\epsilon$ the peak is roughly halved near $\beta_2=0.999$ (right panel), and wherever gradients are smaller than $\epsilon$ Adam becomes gradient descent with step size $\eta/\epsilon$.

\paragraph{How fast the logits drift.} \Cref{prop:spikes:zloss} says that the z-loss removes the flat direction; it does not say whether $\alpha=10^{-4}$ is enough. In a unigram model, a softmax over $K=1000$ classes with Zipf frequencies and logits as the only parameters, trained from zero with minibatches of $64$ labels, stochastic gradient descent keeps the mean logit fixed exactly, as (ii) requires. Adam with $\beta_2=0.95$ does not: $\log Z$ falls steadily, by $0.12$ step sizes per step, while the cross-entropy stays within $0.14$ of its minimum (script \texttt{10-spikes-drift.py}). The drift is a minibatch effect of coordinatewise normalization: it vanishes in full batch, it falls with the batch size ($0.21$, $0.12$, $0.05$, $0.02$, $0.009$ step sizes per step at batches of $16$ to $4096$) and with $\beta_2$ ($0.02$ at $0.99$, $0.001$ at $0.999$), and it comes from the skewed distribution of a rare class's gradient, small and positive in most batches and large and negative in the few that contain it. A z-loss of $10^{-4}$ leaves the drift unchanged; a z-loss of $0.1$ stops it at the price of an excess cross-entropy of $0.38$, and decoupled weight decay of $0.1$ stops it at a price of $0.03$. In this model, then, the coefficient used in practice is far too weak to pin $\log Z$, and weight decay does the pinning; why $10^{-4}$ suffices in large runs is not explained by this mechanism.

\section{What the theory misses}\label{sec:spikes:misses}

\paragraph{What was idealized away.} The main theorems have exact gradients. With noise, a fixed threshold biases the clipped gradient (\cref{ex:spikes:bias}), and \cref{thm:spikes:heavy} needs a threshold that grows with the length of the run and only $L$-smoothness; a theorem that combines heavy tails with $(L_0,L_1)$-smoothness and a fixed practical threshold is more delicate. The constants $L_0$ and $L_1$ are global and fixed, whereas the curvature of a network changes during training, and the assumption of two continuous derivatives excludes ReLU networks. The lower bound is a statement about a class of functions, not about a given network. The burst proposition assumes an exactly constant loud gradient after an exactly silent period; real blocks are rarely silent, and \eqref{eq:spikes:burstbounds} quantifies what a residual state does. The propositions on logits are about invariances and bounds, not dynamics. Finite precision appears nowhere, although it is a candidate for the step from a drifted logit to a divergence.

\paragraph{Where it breaks.} (1) \Cref{ass:spikes:L0L1} fails globally for networks with rescaling symmetries. In the product model $\frac12(uv-1)^2$ of volume 2 of this series, \emph{The Hydrogen Atoms of Deep Learning}, Chapter 15, the minimizers $(u,1/u)$ have zero gradient and sharpness $u^2+u^{-2}$, so no constants $L_0,L_1$ work (\cref{ex:spikes:product}); the assumption holds only on bounded sets, with an $L_0$ that grows with the bound, and then the guarantee depends on how far the weights travel. \Cref{obs:spikes:L0L1} is a correlation along trajectories, not a derivation. (2) Heavy-tailed noise is not the main reason that adaptive methods train transformers faster: \textcite{KunstnerEtAl2023} found that the gap between SGD and Adam persists, and even grows, in full batch, where there is no noise, while sign descent with momentum closes most of it (\cref{ch:transformers}). The case for clipping via \cref{thm:spikes:heavy} is therefore weaker than the case via \cref{thm:spikes:clipped}. (3) Real spikes have interacting causes. In PaLM, neither the batches nor the parameter state alone produced a spike (\cref{obs:spikes:palm}); in OPT, divergences coincided with a collapse of the loss scale of mixed precision \parencite{ZhangEtAl2022OPT}; \textcite{TakaseEtAl2023} (unrefereed) relate spikes to the norms of the Jacobians of the layers and remove them with changes to initialization and to the scale of embeddings, a mechanism of architecture outside all three propositions. (4) In our unigram model the z-loss coefficient of practice does not pin the logits, so the observed benefit of $10^{-4}$ in large runs is unexplained by \cref{prop:spikes:zloss}. (5) MuonClip was needed because, in Kimi K2's experiments, attention logits grew more often with Muon than with AdamW \parencite{KimiTeam2025}; \cref{prop:spikes:qknorm} explains why rescaling removes the symptom, not why an orthogonalized optimizer produces it. (6) That a threshold of $1.0$ is rarely active late in training, except during spikes, is folklore; we know no systematic study.

\paragraph{Repairs and open problems.} The theory of generalized smoothness has grown quickly: clipping with momentum and sharper constants \parencite{ZhangJinFangWang2020}; a symmetric form of \eqref{eq:spikes:L0L1} that needs no second derivatives \parencite{ChenZhouLiangLu2023}; curvature bounded by a general function of the gradient norm, under which gradient descent with a small fixed step converges on sublevel sets \parencite{LiQianTianRakhlinJadbabaie2023}; Adam under such assumptions \parencite{LiRakhlinJadbabaie2023}; generalized sign descent \parencite{CrawshawEtAl2022}; and sharper complexities for gradient methods in the nonconvex and convex cases \parencite{VankovEtAl2025,GorbunovEtAl2025}. For heavy tails there are accelerated and high-probability versions of clipping \parencite{GorbunovDanilovaGasnikov2020} and an analysis of the bias of fixed thresholds \parencite{KoloskovaHendrikxStich2023}. Against the burst, the repairs act on the update: Adafactor-style clipping of the update by its root mean square \parencite{ShazeerStern2018}, adopted for AdamW as StableAdamW \parencite{WortsmanEtAl2023}; spike-aware variants of Adam that clip gradient entries which are large relative to the second-moment estimate and reset the moments \parencite{HuangEtAl2025SPAM}; and Adam-atan2, which replaces $m/(\sqrt v+\epsilon)$ by a bounded function without $\epsilon$ \parencite{EverettEtAl2024}. Against the logits, z-loss, QK-norm, QK-clip, and reparametrizations that control the spectral norm of the attention logits \parencite{ZhaiEtAl2023}. Open: a solvable model in which loss spikes arise and recur without being put in by hand (\cref{ex:spikes:spike}); a derivation of $(L_0,L_1)$-type bounds for a transformer, with constants that depend on the gains of its normalizations; a principled clipping threshold (\cref{ex:spikes:threshold}); and why orthogonalized updates grow attention logits.

\section{Back in the wild}\label{sec:spikes:back}

\fieldnotes{Clip the global gradient norm at $1.0$; $\beta_2=0.95$; z-loss $10^{-4}(\log Z)^2$; QK-norm; MuonClip; after a spike, restart and skip batches.}
{$(L_0,L_1)$-smoothness $\norm{\nabla^2\Loss}\le L_0+L_1\norm{\nabla\Loss}_*$; Adam's moments after a quiet period; the flat direction of the softmax; the bilinear attention logit.}
{A sharp local smoothness lemma. Clipped steepest descent: at most $\max\{4L_0\Delta_0/\eps^2,6L_1\Delta_0/\eps\}$ steps with $\norm g_*>\eps$. Every fixed step: $\gtrsim L_1M\Delta_0/(\eps^2\log(1+L_1M/L_0))$ on some function of the class. Normalized steps are clipped by construction. Clipped SGD at $T^{-(p-1)/(3p-2)}$ with $p$th moments. Exactly: the burst $a_k$, the pinning by the z-loss, the bound of QK-norm by its gains.}
{GD's safe step falls like $\log M/M$, clipped descent's does not; clip where $L_1\norm{\nabla\Loss}_*\approx L_0$; Adam's burst peaks near $0.21/\sqrt{1-\beta_2}$ learning rates about $12$ steps after a block wakes up, six times smaller at $\beta_2=0.95$ than at $0.999$; global clipping does not stop a block's burst.}
{When spikes happen; interacting causes (data, state, precision); $(L_0,L_1)$-smoothness fails under rescaling symmetries; why $10^{-4}$ and $1.0$.}
{Bound the length of every step, by clipping or by normalized updates (proved in the idealization). Prefer a short $\beta_2$ or update clipping when blocks can wake up (suggested). Use z-loss, weight decay and QK-norm to remove directions of growth that the loss cannot see (proved as mechanisms; their effect on spikes is observed). The values $1.0$ and $10^{-4}$ are folklore.}

The theory of this chapter says one thing with confidence: a step whose length is bounded is safe on landscapes whose curvature grows with the gradient, and a fixed step size is not. That is why clipping works, why normalized updates (\cref{ch:steepest,ch:sign,ch:spectral}) do not need it, and why Adam, whose step is bounded by Chapter~2's constant, is less exposed to cliffs than gradient descent; in front of Adam, clipping acts on the moments rather than on the step. It does not tell a practitioner where to put the threshold beyond ``where the curvature starts to grow'', and it predicts no spike. What it adds to the folklore of spikes is a set of exact mechanisms, each of which makes a step too long: an optimizer state that lags its gradient, a direction the loss cannot see, a scale that enters the logits quadratically. Each defence in the practice box removes one of them, and none removes all.

The chapter closes Part III. The stability threshold $2/\eta$ of \cref{ch:eos} bounds the curvature that a fixed step can stand; warmup (\cref{ch:warmup}) grows the step as the curvature falls with the loss; clipping shrinks the step when the curvature rises with the gradient. The loss-dependent curvature of \cref{ch:warmup} and the gradient-dependent curvature of this chapter are the same idea read two ways, and both hold for exponential-type losses. Weight decay (\cref{ch:decay}) is a further defence of the kind studied here: it pins the flat direction of the logits and bounds the scales that enter the attention logits. And the heavy tails of \cref{thm:spikes:heavy} return in \cref{ch:transformers}, where they turn out not to be the main reason that Adam beats gradient descent on transformers.

\begin{exercises}
\item \consolidation (Constants.) (a) Show that $x^4$ is $(L_0,L_1)$-smooth if and only if $L_0L_1^2\ge16$. (b) Show that $e^x$ is $(0,1)$-smooth and $L$-smooth for no $L$. (c) Show that $\log(1+e^{-x})$ is $(\frac14,0)$-smooth and $(0,1)$-smooth. (d) Show that $\sum_ia_i(\cosh\theta_i-1)$ with $a_i>0$ is $(\max_ia_i,1)$-smooth in the Euclidean norm and $(\sum_ia_i,1)$-smooth in $\ell_\infty$. (e) Show that $\frac12\theta^\top H\theta$ with $H\ne0$ is $(L_0,L_1)$-smooth only if $L_0\ge\norm H$. \hint{in (a), minimize $L_0+4L_1x^3-12x^2$ over $x\ge0$; in (c), $\sigma(x)(1-\sigma(x))\le\min\{\frac14,1-\sigma(x)\}$ for the logistic function $\sigma$.}\label{ex:spikes:constants}

\item \consolidation (The lemma is sharp.) (a) Verify the equality case of \cref{lem:spikes:local} for the cliff. (b) Deduce that neither $\varphi_1(c)$ nor $\varphi_2(c)$ in \eqref{eq:spikes:local} can be replaced by a smaller constant. (c) Show that for $\cosh x-1$, which is $(1,1)$-smooth, the bound \eqref{eq:spikes:gradgrowth} is not attained at $\theta=0$, and find the ratio of the two sides as $u\to\infty$.\label{ex:spikes:sharp}

\item \consolidation (Clipping is a trust region.) (a) Show that the clipped step $-\min\{\eta,\gamma/\norm g_2\}g$ minimizes $\ip{g,u}+\norm u_2^2/(2\eta)$ over the ball $\norm u_2\le\gamma$. (b) Show that clipping by value, $u_i=-\sign(g_i)\min\{\eta\abs{g_i},\gamma\}$, minimizes the same function over the box $\norm u_\infty\le\gamma$. (c) Show that clipping by value can change the direction of the step, and global clipping cannot.\label{ex:spikes:trust}

\item \consolidation (The burst.) (a) Show that if the moments were bias-corrected with the age of the coordinate, $\hat m_t=m_t/(1-\beta_1^k)$ and $\hat v_t=v_t/(1-\beta_2^k)$, the step in \cref{prop:spikes:burst} with $\delta=0$ and $\epsilon=0$ would be $\sign(\bar g)$ at every loud step. (b) For $\beta_1=0.9$, show that $(1-\beta_1^k)/\sqrt k$ is largest at $k=12$, and deduce $\max_ka_k\ge6.55$ for $\beta_2=0.999$. (c) For RMSProp, $\beta_1=0$, show that the largest $a_k$ is $a_1=1/\sqrt{1-\beta_2}$. (d) Show that for $\beta_2\to1$ with $\beta_1$ fixed, $\sqrt{1-\beta_2}\,\max_ka_k$ converges to $\max_{k\ge1}(1-\beta_1^k)/\sqrt k$ and $\sqrt{1-\beta_2}\,C(\beta_1,\beta_2)$ to $\sqrt{(1-\beta_1)/(1+\beta_1)}$; for $\beta_1=0.9$ the ratio is $0.903$.\label{ex:spikes:burst}

\item \consolidation (The z-loss.) (a) Show that $\nabla(\log Z)^2=2(\log Z)\pi$ and $\nabla^2(\log Z)^2=2\pi\pi^\top+2(\log Z)(\diag(\pi)-\pi\pi^\top)$. (b) Show that gradient descent on the bias $b$ of a softmax with the loss $\ell+\alpha(\log Z)^2$ changes $\ip{\mathbf1,b}$ by exactly $-2\alpha\eta\log Z$ per step. (c) With decoupled weight decay $\lambda$ instead of the z-loss, show that $\ip{\mathbf1,b}$ is multiplied by $1-\eta\lambda$ per step of gradient descent.\label{ex:spikes:zgrad}

\item \variation (No $(L_0,L_1)$ for the product model.) Let $\Loss(u,v)=\frac12(uv-1)^2$. (a) Show that no constants $L_0,L_1$ make $\Loss$ $(L_0,L_1)$-smooth in the Euclidean norm. (b) Show that $\norm{\nabla^2\Loss(u,v)}_{\op}\le2R^2+1$ on the disk $u^2+v^2\le R^2$. (c) Explain what (b) means for \cref{thm:spikes:clipped} and \cref{cor:spikes:gd} on this model.\label{ex:spikes:product}

\item \variation (Clipping with noise is biased.) Let $g=\mu+\zeta$ with $\mu>0$ and $\zeta=\pm s$ with probability $\frac12$ each, $s>\mu$. (a) Compute $\E\spikesclip_\tau(g)$ for every $\tau>0$, and show that it is $0$ when $\tau\le s-\mu$ and lies in $[0,\mu]$ always. (b) Deduce that clipped SGD with a fixed threshold can stall at a point where the gradient is not zero, and compare with the bias term of \eqref{eq:spikes:heavy}.\label{ex:spikes:bias}

\item \variation (A loud jump.) Suppose that the gradient of a coordinate has been $\delta'\ne0$ for so long that $m=\delta'$ and $v=\delta'^2$ (take $\epsilon=0$), and then jumps to $F\delta'$ with $F>1$. (a) Show that the $k$th step after the jump is $\bigl((1-\beta_1^k)+\beta_1^k/F\bigr)/\sqrt{(1-\beta_2^k)+\beta_2^k/F^2}$, which tends to $a_k$ as $F\to\infty$. (b) Show that the first step after the jump exceeds the steady-state value $1$ for every $F>1$ if $(1-\beta_1)^2\ge1-\beta_2$, and only for $F<(\beta_2-\beta_1^2)/((1-\beta_2)-(1-\beta_1)^2)$ otherwise; evaluate for $(0.9,0.95)$. (c) For $(0.9,0.999)$ find the jump factor at which the first step equals $2$.\label{ex:spikes:jump}

\item \variation (Orthogonalized updates keep the logits centred.) (a) Show that every Newton--Schulz iteration $X\mapsto aX+b(XX^\top)X+c(XX^\top)^2X$ of \cref{ch:spectral} preserves $\mathbf1^\top X=0$, so that Muon, applied to an output layer, leaves $\mathbf1^\top W$ unchanged. (b) Find a gradient $G\in\R^{3\times1}$ with $\mathbf1^\top G=0$ and $\mathbf1^\top\sign(G)\ne0$. (c) Show that the mean logit $\frac1K\mathbf1^\top Wh$ can still change during training when $\mathbf1^\top W$ is fixed, and say through which parameters.\label{ex:spikes:muonmean}

\item \simulation (The band of clipped step sizes.) Rerun \texttt{10-spikes-separation.py} with $\gamma\in\{\frac16,\frac13,\frac23\}$ and $\eta=1/(2L_0)$. Check that the largest convergent step size of clipped descent lies in $[L_1\gamma/(L_0(e^{L_1\gamma/2}-1)),2/L_0]$ for every $M$, that the band shrinks to $2/L_0$ as $\gamma\to0$, and that the number of steps down the cliff grows like $1/\gamma$. Explain why the two-cycle $\pm\gamma/2$ exists exactly when the lower edge is exceeded.\label{ex:spikes:simsep}

\item \simulation (Heavy tails.) Run SGD and clipped SGD on $\Loss(\theta)=\frac12\norm\theta_2^2$ in $\R^{10}$ with noise whose coordinates are independent and symmetric with $\Pr[\abs{\zeta_i}>x]=(1+x)^{-3/2}$, so that $p$th moments are finite exactly for $p<\frac32$. (a) Observe that SGD with a constant step makes rare huge excursions. (b) With the schedule of \cref{thm:spikes:heavy} for $p=1.4$, measure the average of $\norm{\nabla\Loss(\theta_t)}_2^2$ for $T=10^3,\dots,10^6$ and fit a slope in $\log T$; compare it with $-2(p-1)/(3p-2)=-0.36$ (in our runs, about $-0.32$). (c) Explain why the bias term of \eqref{eq:spikes:heavy} is much smaller here than its bound, and which term sets the slope.\label{ex:spikes:simheavy}

\item \simulation (Logit drift.) Run \texttt{10-spikes-drift.py}. (a) Check that SGD keeps the mean logit fixed to rounding error, and that Adam's drift vanishes in full batch. (b) Explain the sign of the drift from the distribution of the gradient of a rare class, small and positive in most batches and large and negative in a few, and the coordinatewise normalization of Adam. (c) Find, by bisection on $\alpha$, the smallest z-loss coefficient that stops the drift at batch size $64$, and compare it with $10^{-4}$.\label{ex:spikes:simdrift}

\item \research (A solvable spike.) Find a low-dimensional model in which Adam with fixed hyperparameters produces recurring loss spikes without a switch put in by hand, and compute their period and height as functions of $\eta$ and $\beta_2$. The burst of \cref{prop:spikes:burst} and the stability threshold of Adam in \cref{ch:eos} are natural ingredients; what is open is a model in which the quiet period itself is produced by the dynamics.\label{ex:spikes:spike}

\item \research (A threshold from data.) Estimate $L_0$ and $L_1$ along a training run by regressing a local smoothness estimate on the gradient norm, as \textcite{ZhangHeSraJadbabaie2020} did, and compare $2L_0/(3L_1)$ with the threshold $1.0$. Is a threshold set from such an estimate better than $1.0$, and how should it change with the size of the model? As far as we know, both questions are open.\label{ex:spikes:threshold}
\end{exercises}

\notesandrefs
Gradient clipping entered deep learning with recurrent language models \parencite{Mikolov2012}, and \textcite{PascanuMikolovBengio2013} gave the norm-clipping rule and the picture of cliffs, building on the analysis of exploding and vanishing gradients by \textcite{BengioSimardFrasconi1994}; the textbook of \textcite{GoodfellowBengioCourville2016} gives both. The rule passed into transformer recipes with a threshold of $1.0$ on the global norm \parencite{BrownEtAl2020,TouvronEtAl2023,ChowdheryEtAl2023}. Adam's step, bounded by construction \parencite{KingmaBa2015}, is clipped before it is computed: the clipped gradient enters its moments.

\Textcite{ZhangHeSraJadbabaie2020} introduced $(L_0,L_1)$-smoothness, measured it along trajectories, and proved the first separation between clipped and fixed-step gradient descent; our proofs follow their strategy, with the exact exponential bound of \cref{lem:spikes:local} in place of its linearization, a maximum in place of a sum in \cref{thm:spikes:clipped}, and the cliff and the plateau as the two functions of \cref{thm:spikes:lower}. The extensions are listed in \cref{sec:spikes:misses}. Heavy-tailed gradient noise was measured by \textcite{SimsekliSagunGurbuzbalaban2019}; \textcite{ZhangEtAl2020Attention} proved the rate of \cref{thm:spikes:heavy} with a matching lower bound and proposed a coordinatewise adaptive clipping; \textcite{KunstnerEtAl2023} showed that noise is not what separates SGD from Adam on transformers.

The record of spikes in large runs is mostly in technical reports: PaLM's restarts and skipped batches and its z-loss \parencite{ChowdheryEtAl2023}, OPT's restarts \parencite{ZhangEtAl2022OPT}, the attention-logit divergence of ViT-22B and its cure by QK-norm \parencite{DehghaniEtAl2023}, and Kimi K2's QK-clip \parencite{KimiTeam2025}. \Textcite{WortsmanEtAl2024} reproduced the two logit instabilities in small models and introduced the sensitivity of the final loss to the learning rate as a measure of stability; \textcite{WortsmanEtAl2023} tied spikes to an out-of-date second moment; \textcite{MolybogEtAl2023} proposed a theory of Adam's instability in large runs; \textcite{ZhaiEtAl2023} tied instability to the collapse of attention entropy. We know no reference for the burst formula \eqref{eq:spikes:burst} beyond its first step, which is the quiet-then-loud step of \cref{ch:adaptive}, nor for the drift of logits under Adam in the unigram model of \cref{sec:spikes:reading}; both are elementary.

\paragraph{Translating the literature.} \Textcite{ZhangHeSraJadbabaie2020} write $f$ and $x$ for our $\Loss$ and $\theta$, and clipped gradient descent as $x_{k+1}=x_k-h_c\nabla f(x_k)$ with $h_c=\min\{\eta_c,\gamma\eta_c/\norm{\nabla f(x_k)}\}$: their $\eta_c$ is our $\eta$, their $\gamma$ is a threshold on the gradient norm, our $\gamma/\eta$, and their $\gamma\eta_c$ is our $\gamma$; their $M$ is ours. \Textcite{ZhangEtAl2020Attention} write $\alpha$ for our moment order $p$ and $G$ for our $\Gamma$. PaLM's z-loss is our $\alpha(\log Z)^2$ with $\alpha=10^{-4}$. Kimi K2's report writes $\tau$ for our logit threshold $\bar s$ and $S_{\max}$ for the largest logit of a head, and multiplies the product of the query and key weights by $\min\{1,\tau/S_{\max}\}$, which is our $\rho^2$. Practitioners' \texttt{max\_norm} is our $\tau=\gamma/\eta$.

\printchapterbib
\end{refsection}
